\section{Introducción}

El problema del tiempo en los sistemas complejos ha recibido aportes decisivos de la termodinámica de los procesos irreversibles y de la teoría del caos. Prigogine mostró que en sistemas alejados del equilibrio surgen estructuras disipativas que introducen una flecha temporal objetiva, irreductible a la reversibilidad de las leyes fundamentales \citep{prigogine_being_to_becoming_1980,prigogine_stengers_order_out_chaos}. Lorenz estableció que atractores deterministas pueden presentar sensibilidad extrema a las condiciones iniciales, de modo que trayectorias inicialmente cercanas divergen de manera exponencial dentro de un conjunto acotado \citep{lorenz_deterministic_nonperiodic_1963}. Feigenbaum demostró que la ruta al caos en una amplia clase de sistemas no lineales obedece a regularidades universales gobernadas por constantes independientes de los detalles microscópicos del modelo \citep{feigenbaum_quantitative_universality_1978}.

Estas perspectivas convergen en que la persistencia en la dinámica de sistemas acoplados no se distribuye de forma uniforme. En regímenes caóticos —como el mapa logístico acoplado con $r=3.8$ y el atractor de Lorenz— la intensificación local de la persistencia, medida mediante hiper-persistencia y trapping estructural, no se traduce automáticamente en reorganización global del sistema. La transición depende de la antisincronización entre componentes, entendida como la desviación estándar de las escalas de persistencia local $\tau_s$. Así, el presente se configura como un fenómeno estratificado y fragmentado, organizado en tres capas jerárquicas cuya integración está regulada por la antisincronización. Esta última actúa como regulador de un efecto estructural de filtrado probabilístico bidireccional que preserva la autonomía relativa de los presentes locales e impide que intensificaciones aisladas produzcan consecuencias estructurales a escala del sistema.

El marco distingue el presente local intensificado (Capa 1), la coherencia entre presentes locales (Capa 2, de naturaleza primariamente relacional y transicional, que introduce la propiedad irreducible de alineación de escalas de persistencia) y la reorganización estructural del sistema (Capa 3). La Capa 2 opera como puente ontológico: su reducción de fragmentación constituye condición necesaria para que las intensificaciones locales adquieran peso estructural. La reformulación de la memoria como retención activa sustentada en trapping e hiper-persistencia completa la caracterización del tiempo extramental discreto.